% Curve-generation runtime — \input{} this file.
\begin{table*}[H]
\centering

\begin{minipage}{16cm}
\centering

\captionsetup{
    justification=centering,
    singlelinecheck=false,
    format=plain
}

\caption{Computational cost of the curve-generation methods.}
\label{tab:timing_results}

\vspace{-0.5em}

\begin{tabular}{p{5.2cm}|c|c|c|c}
\toprule
Method & \makecell{Curve\\models} & \makecell{Training\\total (CPU min)} & \makecell{Training per\\model (CPU s)} & \makecell{Inference per\\curve (ms)} \\
\midrule
ML Step DTW (proposed) & 390 & 246.65 & 37.95 & 22.6 \\
Median per Activity \& Sensor & 390 & \textbf{0.03} & \textbf{$<$0.01} & \textbf{6.6} \\
ML DTW (no steps) & 390 & 870.83 & 133.97 & 20.9 \\
ML (no DTW) & 390 & 869.53 & 133.77 & 16.6 \\
Sensor median & 79 & 0.65 & 0.49 & 6.8 \\
Seq2Seq DTW (aligned) & 390 & 164.83 & 25.36 & 41.5 \\
Seq2Seq (DTW-scored) & 390 & 678.74 & 104.42 & 36.3 \\
\bottomrule
\end{tabular}

\vspace{0.5em}

\parbox{16cm}{%
\footnotesize
Cost of each curve-generation method on the run of Table~\ref{tab:individual_profile_realism}. A curve model is one predictor per (sensor, activity, object); the sensor-median baseline fits one per sensor, so the per-model and per-curve columns are the comparable ones. Training is summed over the 16 workers of the training pool, inference is wall clock over the test curves of a single process. Ratios, not absolute deployment figures. \textbf{Bold} = cheapest per column.
}

\end{minipage}

\end{table*}
